\section{Relative reduction through a specified observation}
\label{sec:relative}

We first give the spectral argument used to pass from geometric estimates
to the complete low space. Its hypotheses concern trial functions,
observed energies and low-energy mass. They do not assume the desired
eigenvalue asymptotics or an identification of the low eigenspace.

\subsection{Forms and geometric estimates}

Let $p$ range over a parameter space with a specified limit. Auxiliary
parameters may vary without converging, provided the following estimates
are uniform. For each $p$, let $\cH_p$ be a real Hilbert space and
$q_p$ a densely defined nonnegative closed quadratic form with compactly
embedded form domain. Let $L_p$ be its associated self-adjoint operator,
and order its eigenvalues, including zeros, as
$\lambda_0(p)\le\lambda_1(p)\le\cdots$.

Independently of this spectrum, suppose that geometry supplies a real
space $V_p$ of dimension $1\le n_p\le N$, a positive definite mass $M_p$,
and a nonnegative energy $E_p$. We use the same letters for their
polarizations. Write
\begin{equation}\label{eq:model-largest}
 \sigma_p=\sup_{x\ne0}\frac{E_p(x)}{M_p(x)}\longrightarrow0,\qquad
 \mu_0(p)\le\cdots\le\mu_{n_p-1}(p)=\sigma_p
\end{equation}
for the largest and all the generalized eigenvalues of $(E_p,M_p)$.
There are three assumptions.

\smallskip\noindent
\emph{Linear trials.}
There is a linear map $T_p:V_p\to\Dom q_p$ such that
\begin{equation}\label{eq:trial-assumption}
 \|T_px\|^2\ge(1-a_p)M_p(x),\qquad
 q_p(T_px)\le(1+a_p)E_p(x),\qquad a_p\to0.
\end{equation}
Thus $T_p$ is injective for sufficiently small parameters. A trial
construction depending on a fixed auxiliary cutoff is enough if, for
each error tolerance, one proves a uniform parameter threshold and one
linear trial map on that domain. A diagonal choice of cutoffs gives
\eqref{eq:trial-assumption}; smooth dependence of that choice is not
needed.

\smallskip\noindent
\emph{Weighted observation.}
A specified linear map $\cA_p:\Dom q_p\to V_p$ satisfies
\begin{equation}\label{eq:observation-assumption}
 E_p(\cA_pu)\le(1+b_p)q_p(u),\qquad b_p\to0.
\end{equation}
The error multiplies the entire energy. An additive $o(\sigma_p)$
bound is insufficient.

\smallskip\noindent
\emph{Low-energy mass capture.}
If $S$ is spanned by at most $n_p+1$ orthonormal eigenvectors with
eigenvalue at most $s$, then
\begin{equation}\label{eq:capture-assumption}
 \bigl|M_p(\cA_pu)-\|u\|^2\bigr|
       \le\omega(p,s)\|u\|^2\quad(u\in S),\qquad
 \lim_{s\downarrow0}\limsup_{p\to0}\omega(p,s)=0.
\end{equation}
We take $\omega$ nonnegative and nondecreasing in $s$. A bound for
every form-domain vector of Rayleigh quotient at most $s$ is stronger
than required. In the geometric applications, this estimate includes
the analysis that excludes additional concentration modes.

No positive lower bound is imposed on the nonzero ratios
$\mu_j/\sigma_p$. The kernels of $E_p$ and $L_p$ may change dimension.
The three assumptions must hold uniformly through the asserted changes.

\begin{theorem}\label{thm:relative-reduction}
Under \eqref{eq:model-largest}--\eqref{eq:capture-assumption}, there
are a fixed $\gamma>0$ and $e_p\to0$ such that the spectral subspace
$S_p$ of $L_p$ below $\gamma$ has dimension $n_p$, and
\begin{equation}\label{eq:abstract-rates}
 (1-e_p)\mu_j(p)\le\lambda_j(p)\le(1+e_p)\mu_j(p)
       \quad(0\le j<n_p),\qquad
 \lambda_{n_p}(p)\ge\gamma.
\end{equation}
The map $\cA_p:S_p\to V_p$ is an isomorphism and maps
$\ker L_p$ onto $\ker E_p$. Define
\begin{equation}\label{eq:abstract-observed-forms}
 N_p(x,y)=\langle\cA_p^{-1}x,\cA_p^{-1}y\rangle_{\cH_p},\qquad
 Q_p(x,y)=q_p(\cA_p^{-1}x,\cA_p^{-1}y).
\end{equation}
Then, as forms on the entire $V_p$,
\begin{equation}\label{eq:abstract-relative-forms}
 (1-e_p)M_p\le N_p\le(1+e_p)M_p,\qquad
 (1-e_p)E_p\le Q_p\le(1+e_p)E_p.
\end{equation}
The error is controlled by the dimension bound and the three hypothesis
errors. It contains no condition number of the positive part of $E_p$.
\end{theorem}

\begin{proof}
Put $c_p=(1+a_p)/(1-a_p)$. Applying min--max to the linear trial
map gives
\begin{equation}\label{eq:abstract-upper}
 \lambda_j\le c_p\mu_j\quad(j<n_p),\qquad
 \lambda_{n_p-1}\le c_p\sigma_p.
\end{equation}
If there were no uniform complementary gap, a sequence tending to the
limit would have $\lambda_{n_p}\to0$. Choose $s>0$ small enough that
$\limsup_{p\to0}\omega(p,s)<1/2$. Eventually the first $n_p+1$
eigenvectors have eigenvalue below $s$ and $\omega(p,s)<3/4$.
Equation~\eqref{eq:capture-assumption} would make $\cA_p$ injective
on their span, contradicting $\dim V_p=n_p$. This proves the gap.
Together with \eqref{eq:abstract-upper}, it identifies precisely
$n_p$ low eigenvalues.

Let $d_p=\omega(p,c_p\sigma_p)$. For each fixed $s>0$, eventually
$c_p\sigma_p\le s$, and hence
$\limsup_{p\to0}d_p\le\limsup_{p\to0}\omega(p,s)$. Letting $s\downarrow0$
shows $d_p\to0$, also when $\sigma_p=0$. The mass-capture estimate
makes $\cA_p$ injective on $S_p$, hence an isomorphism, and gives
\begin{equation}\label{eq:observed-mass-bounds}
 (1+d_p)^{-1}M_p\le N_p\le(1-d_p)^{-1}M_p.
\end{equation}

The trial map sends $\ker E_p$ injectively into $\ker L_p$. Conversely,
the observation inequality sends $\ker L_p$ to $\ker E_p$, and
mass capture is injective on this actual kernel. The two nullities
are therefore equal and $\cA_p$ identifies the kernels exactly.

For $j<n_p$, the image under $\cA_p$ of the first $j+1$ eigenvectors
has dimension $j+1$. Finite-dimensional min--max gives a nonzero
vector $u$ in that span such that
$E_p(\cA_pu)\ge\mu_jM_p(\cA_pu)$. Thus
\[
 (1+b_p)\lambda_j\|u\|^2
       \ge E_p(\cA_pu)
       \ge\mu_j(1-d_p)\|u\|^2,
\]
and consequently
\begin{equation}\label{eq:abstract-lower}
 \lambda_j\ge\frac{1-d_p}{1+b_p}\mu_j.
\end{equation}
This proves the relative eigenvalue comparison. The mass statement
already follows from \eqref{eq:observed-mass-bounds}.

For the stronger energy statement,
\eqref{eq:observation-assumption} gives
$E_p\le(1+b_p)Q_p$. Their kernels agree by the preceding argument;
denote this common kernel by $K_p$. Both energies descend to positive
forms on $V_p/K_p$. On this quotient define the masses
\begin{equation}\label{eq:quotient-masses}
 \overline M_p([x])=\inf_{k\in K_p}M_p(x+k),\qquad
 \overline N_p([x])=\inf_{k\in K_p}N_p(x+k).
\end{equation}
Taking infima preserves \eqref{eq:observed-mass-bounds}. The positive
generalized eigenvalues of the original pairs are exactly the
generalized eigenvalues of their quotient pairs: in each class the
mass minimizer is perpendicular to $K_p$ for the corresponding mass,
and the energy is constant on that class.

Put $k_p=\dim(V_p/K_p)\le N$. If $k_p=0$, both energies vanish and
there is nothing to prove. Otherwise their quotient determinants in
any common basis satisfy
\begin{equation}\label{eq:quotient-determinant}
 \frac{\det\overline Q_p}{\det\overline E_p}
   =\prod_{\lambda_j>0}\frac{\lambda_j}{\mu_j}
           \frac{\det\overline N_p}{\det\overline M_p}.
\end{equation}
In fixed Euclidean coordinates on the quotient, consider
$B_p=\overline E_p^{-1/2}\overline Q_p\overline E_p^{-1/2}$.
Every eigenvalue of $B_p$ is at least $(1+b_p)^{-1}$. The upper
eigenvalue and mass bounds imply
\[
 \det B_p\le c_p^{k_p}(1-d_p)^{-k_p}.
\]
Dividing this product by the lower bounds for the other $k_p-1$
eigenvalues gives
\begin{equation}\label{eq:determinant-upgrade}
 \lambda_{\max}(B_p)
       \le c_p^{k_p}(1-d_p)^{-k_p}(1+b_p)^{k_p-1}.
\end{equation}
Both bounds tend to one uniformly for $k_p\le N$, independently of
the smallest positive energy. Hence the quotient energies are
relatively equivalent. Their identical kernels lift this equivalence
to all of $V_p$, proving \eqref{eq:abstract-relative-forms}.
\end{proof}

\begin{corollary}\label{cor:observed-matrices}
Let $u_0,\ldots,u_{n_p-1}$ be an orthonormal eigenbasis of $S_p$,
let $C_p$ have columns $\cA_pu_j$, and put
$\Lambda_p=\operatorname{diag}(\lambda_0,\ldots,\lambda_{n_p-1})$.
After enlarging $e_p$ by a fixed factor,
\begin{equation}\label{eq:general-observed-matrices}
 \bigl|M_p(C_pz)-|z|^2\bigr|\le e_p|z|^2,\qquad
 \bigl|E_p(C_pz)-z^T\Lambda_pz\bigr|
                            \le e_p z^T\Lambda_pz.
\end{equation}
Each positive-mode cross-error in the energy matrix is bounded by
$e_p\sqrt{\lambda_i\lambda_j}$; every kernel row and column is
exactly zero.
\end{corollary}
\begin{proof}
The definitions give $C_p^TN_pC_p=I$ and $C_p^TQ_pC_p=\Lambda_p$.
Conjugate \eqref{eq:abstract-relative-forms} by $C_p$. On the positive
block, conjugation also by $\Lambda_p^{-1/2}$ bounds the normalized
operator norm and hence every entry. Equality of kernels gives the
remaining rows and columns.
\end{proof}

\subsection{Why the observation matters}

The conclusion is a comparison in the prescribed geometric coordinates.
A near-identity change to a different orthonormal basis need not preserve
relative error at the weakest scale. For example, let
$E_\varepsilon=\operatorname{diag}(\varepsilon^2,\varepsilon)$ and
let $R_\theta$ be a Euclidean rotation with $\theta\to0$. Then
\[
 \|R_\theta^TE_\varepsilon R_\theta-E_\varepsilon\|
       =o(\varepsilon),
\]
whereas the ratio of the first diagonal entries is
$\cos^2\theta+\varepsilon^{-1}\sin^2\theta$. For
$\theta=\varepsilon^{1/4}$ this ratio diverges. Rotating a zero
eigendirection similarly changes the kernel in fixed coordinates.
An additive estimate at the largest small eigenvalue scale therefore does not
give \eqref{eq:abstract-relative-forms}.

This does not prevent a change of coordinates in the same geometric
construction. For an invertible $S_p^f:V'_p\to V_p$, take
\[
 M'_p=(S_p^f)^*M_p,\quad E'_p=(S_p^f)^*E_p,\quad
 T'_p=T_pS_p^f,\quad \cA'_p=(S_p^f)^{-1}\cA_p.
\]
All hypotheses and conclusions transform by exact congruence. There is
no Euclidean condition-number assumption on $S_p^f$, since all the
forms and maps are transported together.

In the applications, an energy may be written $E=d^*Wd$, where the
adjoint uses a declared reference inner product and $W$ is positive on
the target. Its mass-normalized matrix is
$M^{-1/2}d^*WdM^{-1/2}$. This elementary factorization does not
construct an entire elliptic complex or determine the spectrum in
another degree. Nor does Theorem~\ref{thm:relative-reduction} identify
$Q_p(x)$ with the minimum of $q_p$ over all sections satisfying
$\cA_pu=x$. It is the energy on the actual low spectral subspace.
No exact identity $\cA_pT_p=I$ is assumed.
